\scanpage{178}
\mspage{127}
\diagram{\begin{tikzcd}[ampersand replacement=\&,column sep=small,row sep=small]
\& \& X_1 \& \& \& X'_1\\
f_{\mathscr X}=\int f:\quad\mathscr X=\int \& X_0 \arrow[ru]\arrow[rd] \& \& {\longrightarrow}\quad\mathscr X'=\int \& X'_0 \arrow[ru]\arrow[rd] \& \\
\& \& X_2 \& \& \& X'_2\rlap{,}
\end{tikzcd}}
s'insérant dans un \deleted{cube} \inserted{diagramme cubique} (cf diagramme précédent complété par des pointillés). Considérons les conditions suivantes
\[
\begin{aligned}
\text{a)}\quad & f_0\in W\\
\text{b)}\quad & f_1\text{ et }f_2\text{ sont }\in W\\
\text{c)}\quad & f_{\mathscr X}=\int(f_0,f_1,f_2)\quad\in W
\end{aligned}
\]
Supposons l'axiome $W(7)$ (\textquotedblleft Lemme des 5\textquotedblright, cf page 76) satisfait. Alors deux \deleted{\unsure{quel}} parmi les conditions $a,b,c$ impliquent la troisième :
\[
\left\{\begin{aligned}
a\mathbin{\&} b&\Longrightarrow c\\
b\mathbin{\&} c&\Longrightarrow a\\
c\mathbin{\&} a&\Longrightarrow b .
\end{aligned}\right.
\]

En effet, on a un hom du \deleted{\unsure{diagra}} \unsure{mme|même} cubique ci-dessus dans le diagramme similaire, où $\deleted{\unsure{\alpha}}X_1,X_2,X'_1,X'_2$ sont remplacés respectivement par les \deleted{\unsure{ouverts}}\uline{\unsure{fermés|souscatég.}} $\,{}_{1/}\mathscr X_{\deleted{\unsure{/1}}},{}_{2/}\mathscr X_{\deleted{\unsure{/2}}}$ de $\mathscr X$, et $\,{}_{1/}\mathscr X'_{\deleted{\unsure{/1}}},{}_{2/}\mathscr X'_{\deleted{\unsure{/2}}}$ de $\mathscr X'$, les deux carrés sup. et inférieur étant les inclusions d'\unsure{ouverts|éléments} de $\mathscr X$ resp. de $\mathscr X'$. Ainsi, les conditions $a),b),c)$ \inserted{sont} équivalentes respectivement aux conditions similaires pour le nouveau diagramme cubique
% archive sheet number: 178
